\documentclass{article}
\usepackage{lmodern}
\usepackage{amsmath}
\usepackage{listings}
\usepackage{fullpage}
\usepackage{parskip}
\usepackage{xcolor}
\lstset{
	basicstyle=\ttfamily,
	columns=fullflexible,
	frame=single,
	breaklines=true,
	postbreak=\mbox{\textcolor{red}{$\hookrightarrow$}\space},
}
\begin{document}
\title{Experiment `p\_6\_2\_a\_swap\_first\_last\_array' Results}
\author{\today}
\date{}
\maketitle
\textbf{Experiment outcome: }SUCCESS
\\\textbf{Bad responses: }0
\\\textbf{Responses containing}~\texttt{assume}~\textbf{: }0
\\\textbf{Responses containing}~\texttt{expect}~\textbf{: }0
\\\textbf{Resolution attempts: }2
\\\textbf{Hard fails (resolution): }0
\\\textbf{Soft fails (resolution): }0
\\\textbf{Verification attempts: }2
\section*{Problem Specification}
\textbf{Problem name: }p\_6\_2\_a\_swap\_first\_last\_array
\\\textbf{Natural language statement: }Write a method to swap the first and last elements in an array.
\\\textbf{Method signature: }p\_6\_2\_a\_swap\_first\_last\_array(inputs: array<int>) returns (result: array<int>)
\subsection*{Ensures}
\begin{itemize}
\item \texttt{result.Length == inputs.Length}
\item \texttt{result[0] == inputs[inputs.Length - 1]}
\item \texttt{result[inputs.Length - 1] == inputs[0]}
\item \texttt{forall i :: 1 <= i < inputs.Length - 1 ==> result[i] == inputs[i]}
\end{itemize}
\subsection*{Requires}
\begin{itemize}
\item \texttt{inputs.Length >= 2}
\end{itemize}
\clearpage
\section*{GenAI interactions}
Below you will find all interactions between the `user' (program) and the `assistant' (OpenAI).
\subsection*{Program $\rightarrow$ GenAI}
\begin{lstlisting}
You are given the following task to perform in Dafny:

Write a method to swap the first and last elements in an array.

The signature should be:

p_6_2_a_swap_first_last_array(inputs: array<int>) returns (result: array<int>)

The method should respect the following contract:

requires inputs.Length >= 2, ensures result.Length == inputs.Length, ensures result[0] == inputs[inputs.Length - 1], ensures result[inputs.Length - 1] == inputs[0], ensures forall i :: 1 <= i < inputs.Length - 1 ==> result[i] == inputs[i]

Produce and show only the Dafny body of this method, including the curly braces that surround it. Do not show the signature nor contract. You must not use 'assume' nor 'expect' anywhere in your code.
\end{lstlisting}
\subsection*{GenAI $\rightarrow$ Program}
\textbf{System fingerprint: }null
\\\textbf{ID: }chatcmpl-Dn0BYb8uF49vGABxwtCnM28sfJMzV
\\\textbf{Model: }gpt-5.5-2026-04-23
\\\textbf{Created at: }1780571156
\\\textbf{Finish reason: }stop
\\\textbf{Usage: }378 tokens in, and 587 tokens out
\begin{lstlisting}
{
  result := new int[inputs.Length](i =>
    if i == 0 then inputs[inputs.Length - 1]
    else if i == inputs.Length - 1 then inputs[0]
    else inputs[i]
  );
}
\end{lstlisting}
\subsection*{Program $\rightarrow$ GenAI}
\begin{lstlisting}
When using dafny verify, the below error is emitted and verify fails:

problem_data/p_6_2_a_swap_first_last_array/program.dfy(9,25): Error: insufficient reads clause to read array element; Consider extracting inputs[inputs.Length - 1] to a local variable before the lambda expression, or adding 'reads inputs' in the enclosing lambda specification for resolution
  |
9 |     if i == 0 then inputs[inputs.Length - 1]
  |                          ^

problem_data/p_6_2_a_swap_first_last_array/program.dfy(10,46): Error: insufficient reads clause to read array element; Consider extracting inputs[0] to a local variable before the lambda expression, or adding 'reads inputs' in the enclosing lambda specification for resolution
   |
10 |     else if i == inputs.Length - 1 then inputs[0]
   |                                               ^

problem_data/p_6_2_a_swap_first_last_array/program.dfy(11,15): Error: index out of range
   |
11 |     else inputs[i]
   |                ^

problem_data/p_6_2_a_swap_first_last_array/program.dfy(11,15): Error: insufficient reads clause to read array element; Consider extracting inputs[i] to a local variable before the lambda expression, or adding 'reads inputs' in the enclosing lambda specification for resolution
   |
11 |     else inputs[i]
   |                ^


Dafny program verifier finished with 1 verified, 4 errors

Correct the error in code by altering only the method body. Produce and show only the Dafny body, including the curly braces that surround it. Do not show the signature nor contract. You must not use 'assume' nor 'expect' anywhere in your code.
\end{lstlisting}
\subsection*{GenAI $\rightarrow$ Program}
\textbf{System fingerprint: }null
\\\textbf{ID: }chatcmpl-Dn0C9ABvC7L1AFVIQcCciiqw1vpyC
\\\textbf{Model: }gpt-5.5-2026-04-23
\\\textbf{Created at: }1780571193
\\\textbf{Finish reason: }stop
\\\textbf{Usage: }1017 tokens in, and 893 tokens out
\begin{lstlisting}
{
  result := new int[inputs.Length];
  result[0] := inputs[inputs.Length - 1];
  result[inputs.Length - 1] := inputs[0];

  var i := 1;
  while i < inputs.Length - 1
    invariant 1 <= i <= inputs.Length - 1
    invariant result.Length == inputs.Length
    invariant result[0] == inputs[inputs.Length - 1]
    invariant result[inputs.Length - 1] == inputs[0]
    invariant forall j :: 1 <= j < i ==> result[j] == inputs[j]
  {
    result[i] := inputs[i];
    i := i + 1;
  }
}
\end{lstlisting}
\section*{Final Program}
\begin{lstlisting}
method p_6_2_a_swap_first_last_array(inputs: array<int>) returns (result: array<int>)
	requires inputs.Length >= 2
	ensures result.Length == inputs.Length
	ensures result[0] == inputs[inputs.Length - 1]
	ensures result[inputs.Length - 1] == inputs[0]
	ensures forall i :: 1 <= i < inputs.Length - 1 ==> result[i] == inputs[i]
{
  result := new int[inputs.Length];
  result[0] := inputs[inputs.Length - 1];
  result[inputs.Length - 1] := inputs[0];

  var i := 1;
  while i < inputs.Length - 1
    invariant 1 <= i <= inputs.Length - 1
    invariant result.Length == inputs.Length
    invariant result[0] == inputs[inputs.Length - 1]
    invariant result[inputs.Length - 1] == inputs[0]
    invariant forall j :: 1 <= j < i ==> result[j] == inputs[j]
  {
    result[i] := inputs[i];
    i := i + 1;
  }
}
\end{lstlisting}
\section*{Total Token Usage}
\textbf{Input tokens: }1395
\\\textbf{Output tokens: }1480
\\\textbf{Reasoning tokens: }1219
\\\textbf{Sum of `total tokens': }2875
\section*{Experiment Timings}
\textbf{Overall Experiment} started at 1780571156024, ended at 1780571228535, lasting 72511ms (72.51 seconds)
\\\textbf{Iteration \#1} started at 1780571156025, ended at 1780571192732, lasting 36707ms (36.71 seconds)
\\\textbf{Iteration \#2} started at 1780571192733, ended at 1780571228534, lasting 35801ms (35.80 seconds)
\end{document}
